\input{preamble}

% https://phys.libretexts.org/Bookshelves/Quantum_Mechanics/Introductory_Quantum_Mechanics_(Fitzpatrick)/04%3A_One-Dimensional_Potentials/4.05%3A_Alpha_Decay

\section*{Alpha decay 3}

In this section we derive $C_3$ from the following half-life data.
Rows with $Z=90$ are uranium and $Z=88$ are thorium.
Recall that $Z$ is atomic number minus two.
\begin{equation*}
\begin{matrix}
\log\tau_{1/2} & A & Z & E
\\[1ex]
6.30 & 228 & 90 & 6.80
\\
14.4 & 230 & 90 & 5.99
\\
21.5 & 232 & 90 & 5.41
\\
29.7 & 234 & 90 & 4.86
\\
34.2 & 236 & 90 & 4.57
\\
39.5 & 238 & 90 & 4.27
\\
7.52 & 226 & 88 & 6.45
\\
17.9 & 228 & 88 & 5.52
\\
28.5 & 230 & 88 & 4.77
\\
40.6 & 232 & 88 & 4.08
\end{matrix}
\end{equation*}

From the previous section we have
\begin{equation*}
\log\tau_{1/2}=2\left(C_1\frac{Z}{\sqrt E}-C_2\sqrt{Zr_1}\right)+C_3\log A+\log(\log2)
\end{equation*}

Hence the regression model
\begin{equation*}
\log\tau_{1/2}-2\left(C_1\frac{Z}{\sqrt E}-C_2\sqrt{Zr_1}\right)-\log(\log2)=C_3\log A
\end{equation*}

By ordinary least squares we obtain
\begin{equation*}
C_3=-10.1
\end{equation*}

The regression formula yields the following predicted values.
\begin{equation*}
\begin{matrix}
& \text{Predicted}
\\
\log\tau_{1/2} & \log\tau_{1/2}
\\[1ex]
6.30 & 5.31
\\
14.4 & 14.0
\\
21.5 & 21.4
\\
29.7 & 29.7
\\
34.2 & 34.5
\\
39.5 & 40.1
\\
7.52 & 6.87
\\
17.9 & 17.8
\\
28.5 & 28.8
\\
40.6 & 41.6
\end{matrix}
\end{equation*}

\href{https://georgeweigt.github.io/examples/alpha-decay-3-demo.html}{Eigenmath script}

\end{document}
